\textbf{COLUNI 2022 -- Questão 30}\par\medskip A Floresta Atlântica é considerada um dos cinco mais importantes \textit{hotspots} de biodiversidade (pontos quentes; de alta biodiversidade) mundiais, da qual resta menos de 8\% da cobertura original. De um total estimado em mais de 1800 espécies de vertebrados terrestres, aproximadamente 660 (36\%) são endêmicos ao bioma (só ocorrem na Mata Atlântica). Dos 305 táxons de animais brasileiros listados como ameaçados pela União Internacional para a Conservação da Natureza (IUCN), 188 ocorrem na Mata Atlântica. Desses, 151 são vertebrados e 116 endêmicos.

\par

PAGLIA, A.P., A. PAESE, L. BEDÊ, M. FONSECA, L.P. PINTO E MACHADO R.B.. 2004. Lacunas de conservação e áreas insubstituíveis para vertebrados ameaçados da Mata Atlântica. Pp. 39-50. \textit{In:} Anais do IV Congresso Brasileiro de Unidades de Conservação. Volume II - Seminários. Fundação o Boticário de Proteção à Natureza e Rede Nacional Pró Unidades de Conservação. Curitiba, PR. (Adaptado).

\par

O texto anterior trata da biodiversidade de vertebrados terrestres na Mata Atlântica. Sobre os vertebrados em geral, é INCORRETO afirmar que:\begin{enumerate}
\item os ovos com casca constituem um importante passo evolutivo na conquista do ambiente terrestre pelos
      répteis.
\item a mandíbula está presente em todos os vertebrados e tem importante papel na captura do alimento e na
      defesa contra predadores.
\item os peixes são representantes importantes dos vertebrados aquáticos e apresentam características
      particulares de adaptação à água, como a respiração branquial.
\item a notocorda é uma estrutura presente na fase embrionária desses organismos, o que os inclui no grupo
      dos Chordata.\end{enumerate}
